\documentclass[12pt]{article}
\usepackage[utf8]{inputenc}
\usepackage[T1]{fontenc}
\usepackage{hyperref}
\title{Middle Mongol Language Resources for K-12 Curriculum Development}
\author{Compiled by Claude Sonnet 5.5 (claude-sonnet-5-5) for Andr\'as Kornai}
\date{}
\begin{document}
\maketitle

\section{Introduction}
Middle Mongol (ISO 639-3 xng; Glottocode midd1351) is the Mongolic language of the Mongol Empire and its successor states. It is attested in manuscripts and inscriptions in several scripts. It has no native speakers today, so any school programme would be a heritage and classical-language programme. It would build on Modern Mongolian and on Written Mongolian. The search found only a few specialised scholarly resources and no school materials. Everything below is scholarly or reference material.

\section{Primary Educational Resources}
No primers, textbooks or school curricula for Middle Mongol turned up in the search. Teachers would have to adapt scholarly editions and the Modern Mongolian materials listed below.

\section{Linguistic and Reference Materials}
\begin{itemize}
\item Overview article on Middle Mongol, including its relation to Modern Mongolian: \url{https://wikipedia.com/wiki/ISO_639:xng}
\item Encyclopedia entry on Middle Mongol: \url{https://encyclopedia.pub/entry/32585}
\item Monumenta Altaica, a page on Poppe's \emph{Mongol'skij slovar' Mukaddimat al-adab}. This is a Mongolian--Turkic lexicon from a 14th-century manuscript, and it contains full sentences: \url{https://altaica.ru/e_mukadd.php}
\item Mirror of the same Monumenta Altaica page: \url{https://altaica.c4.starlingdb.org/e_mukadd.php}
\item Lexilogos Mongolian dictionary page. It lists dictionaries and grammars of Modern and Written Mongolian, which are indirect aids: \url{https://lexilogos.com/english/mongolian_dictionary.htm}
\item \emph{A Modern Mongolian-English Dictionary}, Indiana University Press. This covers Modern Mongolian only: \url{https://iupress.org/9780933070196/a-modern-mongolian-english-dictionary}
\end{itemize}

\section{Religious and Community Resources}
The search found no community organisations or religious resources specific to Middle Mongol. Related literary and historical texts are known to scholarship, but none were located online in this search.

\section{Digital and Audio Resources}
No audio, software or corpus resources specific to Middle Mongol were found in the search. Audio would have to be reconstructed or drawn from Modern Mongolian.

\section{Curriculum Development Context}
Middle Mongol cannot support a curriculum that uses the language as the medium of instruction. A realistic design is a short heritage unit within Mongolian-language or history classes. It could cover the language's history, its scripts, and selected short texts from the Mukaddimat al-adab. Content would need expert review, because the search found no pedagogical grammar. A first step is to find a specialist in Mongolic historical linguistics.
\end{document}